\documentclass[11pt]{article}
\usepackage[localtheorems]{raeez-math-template}
\newtheorem{theorem}{Theorem}[section]
\newtheorem{proposition}[theorem]{Proposition}
\newtheorem{lemma}[theorem]{Lemma}
\newtheorem{definition}[theorem]{Definition}
\newtheorem{conjecture}[theorem]{Conjecture}
\newtheorem{remark}[theorem]{Remark}

\newcommand{\cA}{\mathcal{A}}
\newcommand{\cC}{\mathcal{C}}
\newcommand{\cD}{\mathcal{D}}
\newcommand{\cH}{\mathcal{H}}
\newcommand{\cM}{\mathcal{M}}
\newcommand{\cO}{\mathcal{O}}
\newcommand{\cV}{\mathcal{V}}
\newcommand{\cZ}{\mathcal{Z}}
\newcommand{\fg}{\mathfrak{g}}
\newcommand{\bC}{\mathbb{C}}
\newcommand{\bR}{\mathbb{R}}
\newcommand{\bZ}{\mathbb{Z}}
\newcommand{\bP}{\mathbb{P}}
\newcommand{\CY}{\mathrm{CY}}
\newcommand{\Im}{\operatorname{Im}}
\newcommand{\Vol}{\operatorname{Vol}}
\newcommand{\Jac}{\operatorname{Jac}}
\newcommand{\AJ}{\operatorname{AJ}}
\newcommand{\Tr}{\operatorname{Tr}}
\newcommand{\kch}{\kappa_{\mathrm{ch}}}
\newcommand{\kcat}{\kappa_{\mathrm{cat}}}

\title{The Chiral Volume Conjecture:\\
Abel--Jacobi Periods from Chiral Partition Functions}
\author{Raeez Lorgat}
\date{2026-04-22}

\begin{document}
\maketitle

\begin{abstract}
On a compact Calabi--Yau threefold $X$ containing a holomorphic
curve $C$, the CY-to-chiral functor $\Phi$ (proved equivalence
at $d=3$: Theorem~CY-A$_3$) produces an $E_1$-chiral algebra
$\cA_X$ whose representation theory at level $N$ is a finite-rank
module over the affine Grassmannian of the K3 Yangian transported
along $C$. The \emph{chiral volume conjecture} asserts that the
large-$N$ asymptotics of the chiral partition function $Z_N(\cA_X; C)$
are governed by the Abel--Jacobi period of $C$ against the
holomorphic three-form: concretely,
\[
  \lim_{N \to \infty} \frac{1}{N}\, \log \bigl|Z_N(\cA_X; C)\bigr|
  \;=\; \frac{1}{2\pi}\, \bigl|\AJ_X(C)\bigr|.
\]
We state the conjecture with each object defined at its proper
ambient, derive the factor $1/2\pi$ from semiclassical $6$d
holomorphic Chern--Simons on the total space of
$\Omega^{3,0}_X \to X$, and identify the right-hand side as the
on-shell action of the classical $6$d hCS configuration sourced
by $C$. The statement is conjectural; it is conditional on
CY-A$_3$ (now proved) and on the modular-sum identification of
$Z_N$ with the level-$N$ Verlinde-type generating function.
\end{abstract}

\section{The ambient: CY$_3$, holomorphic $3$-form, and chiral
datum}
\label{sec:ambient}

Fix a smooth compact Calabi--Yau threefold
$(X, \Omega_X)$ with $\Omega_X \in H^{3,0}(X, \bC)$ the holomorphic
volume form, normalised so that
$\int_X \Omega_X \wedge \overline{\Omega_X} = 1$. Let $C \subset X$
be a smooth projective curve, either rigid or embedded in a flat
family.

The cohomology $H^3(X, \bC)$ carries the pure Hodge structure
$H^3 = H^{3,0} \oplus H^{2,1} \oplus H^{1,2} \oplus H^{0,3}$.
The intermediate Jacobian is
\[
  \mathrm{J}^2(X) \;=\; \frac{H^3(X, \bC)}{F^2 H^3(X, \bC) \oplus H^3(X, \bZ)},
\]
a complex torus of dimension $h^{2,1}(X) + 1$ (viewed through the
Weil complex structure).

\begin{definition}[Abel--Jacobi invariant of a curve]
\label{def:aj}
For a null-homologous $1$-cycle $C \in Z^{1}(X)$ (equivalently a
divisor on $C$ of degree zero), Griffiths' Abel--Jacobi map is
\[
  \AJ_X\colon Z^{2}_{\mathrm{hom}}(X) \longrightarrow \mathrm{J}^2(X),
  \qquad
  C \longmapsto \left[ \omega \mapsto \int_{\Gamma} \omega \right],
\]
where $\partial \Gamma = C$ and $\omega$ ranges over $F^2 H^3(X, \bC)$.
The magnitude $|\AJ_X(C)|$ is measured in the Weil metric on
$\mathrm{J}^2(X)$; equivalently,
\[
  \bigl|\AJ_X(C)\bigr|
  \;=\;
  \left| \int_{\Gamma} \Omega_X \right|
  \;+\; O\!\left( \bigl|\int_\Gamma H^{2,1}\bigr| \right),
\]
with the first term dominant because $\Omega_X$ is the unique
$(3,0)$-piece.
\end{definition}

For a curve $C$ that is the boundary of a Lefschetz thimble
$\Gamma_C$ in a family, the leading term
$\int_{\Gamma_C} \Omega_X$ is the \emph{period} of $\Omega_X$ over
$\Gamma_C$. This is the only complex-structure invariant that
enters: it does not depend on the K\"ahler class. The CY$_3$
volume $\int_X \omega^3 / 3!$ (K\"ahler-dependent) is irrelevant
here, which is why the chiral version of Kashaev's conjecture
succeeds where the na\"ive ``volume of $X$'' version fails.

The CY-to-chiral functor $\Phi$ at $d=3$ (Theorem~CY-A$_3$) sends
the CY-category $\mathrm{D}^b(X)$ to an $E_1$-chiral algebra $\cA_X$,
whose representation category $\Rep(\cA_X)$ carries an $E_2$-structure
realised on the Drinfeld centre $Z(\Rep(\cA_X))$. The curve $C$ is
the \emph{factorisation locus}: points of $C$ label insertions of
vertex operators, so $C$ acts as a Ran-space probe into $\cA_X$.

\section{The chiral partition function $Z_N(\cA_X; C)$}
\label{sec:ZN}

\begin{definition}[Level-$N$ chiral partition function]
\label{def:ZN}
Let $\cH_N(\cA_X; C)$ be the space of level-$N$ conformal blocks of
$\cA_X$ on $(C, \text{trivial insertion})$; equivalently, the
$\Sym^N$-equivariant sector of the Fock-type module constructed from
the factorisation $E_1$-structure of $\cA_X$ pulled back along
$C \hookrightarrow X$. Define
\[
  Z_N(\cA_X; C)
  \;:=\; \Tr_{\cH_N(\cA_X; C)} q^{L_0 - c/24}\Big|_{q = \exp(2\pi i \tau_C)},
\]
evaluated at the modulus $\tau_C$ naturally attached to $C$ by the
period of the holomorphic three-form along a thimble
$\Gamma_C$: $\tau_C := \int_{\Gamma_C} \Omega_X$.
\end{definition}

Two features are essential. First, $Z_N$ is intrinsically
chiral-side: the trace is over the chiral Hilbert space, not over
geometric cycles on $X$. Second, the modulus $\tau_C$ is the
Abel--Jacobi period of $C$ under $\Omega_X$ --- the complex
parameter governing the symmetric product structure.

The symmetric-product structure is forced by the $E_1$-chiral
nature of $\cA_X$ at $d = 3$: the factorisation product on
$\cA_X$ along a $1$-parameter family of points in $C$ produces a
graded module over $\Sym^\bullet(\cA_X)$, whose $N$-th graded
piece is $\cH_N$.

\section{Statement of the conjecture}
\label{sec:conj}

\begin{conjecture}[Chiral volume conjecture]
\label{conj:chiral-volume}
Let $X$ be a smooth compact CY$_3$, let $C \subset X$ be a smooth
null-homologous curve, let $\cA_X = \Phi(\mathrm{D}^b(X))$ be the
$E_1$-chiral algebra produced by the CY-to-chiral functor
(CY-A$_3$), and let $Z_N(\cA_X; C)$ be the level-$N$ chiral
partition function of Definition~\ref{def:ZN}. Then
\begin{equation}
  \label{eq:chiral-volume}
  \lim_{N \to \infty} \frac{1}{N}\, \log \bigl|Z_N(\cA_X; C)\bigr|
  \;=\; \frac{1}{2\pi}\, \bigl|\AJ_X(C)\bigr|.
\end{equation}
\end{conjecture}

The right-hand side is the Weil-metric magnitude of the Abel--Jacobi
class in $\mathrm{J}^2(X)$; equivalently, to leading order,
$(2\pi)^{-1} |\!\int_{\Gamma_C} \Omega_X|$.

The statement is conditional on CY-A$_3$ (proved at the
$(\infty,1)$-categorical level at $d=3$, inscribed in
\texttt{chapters/theory/cy\_to\_chiral.tex} and in
\texttt{working\_notes.tex} \S\ref{subsec:chiral-volume-conjecture})
and on the identification of $Z_N$ with the generating function of
level-$N$ Verlinde-type dimensions --- the latter is the actual
mathematical content to be proved.

\section{Attack cycle 1 --- the coefficient $1/(2\pi)$}
\label{sec:attack1}

\emph{Attack.} Kashaev's conjecture reads
$\lim_N \frac{2\pi}{N} \log |J_N(K; e^{2\pi i /N})| = \Vol(S^3 \setminus K)$,
with coefficient $2\pi$, not $1/(2\pi)$. Why is the chiral analogue
inverted?

\emph{Heal.} The coefficients are consistent, not inverted. In
Kashaev's formula, $J_N$ is evaluated at the root of unity
$e^{2\pi i /N}$ where the \emph{growth rate} scales as $N/(2\pi)$
times the volume, yielding $(2\pi/N) \log|J_N| \to \Vol$. In
equation~\eqref{eq:chiral-volume}, the trace is over a level-$N$
module with modulus $\tau_C = \int_{\Gamma_C} \Omega_X$ fixed (not
$N$-dependent), and the large-$N$ limit is of
$\log Z_N/N \to |\AJ_X(C)|/(2\pi)$. The two identities rearrange to
\[
  \frac{1}{N} \log |Z_N|  \;\sim\;  \frac{1}{2\pi}|\AJ_X(C)|,
  \qquad
  \frac{2\pi}{N} \log |J_N| \;\sim\; \Vol(S^3 \setminus K).
\]
Multiplying the first by $2\pi$ gives
$(2\pi/N) \log |Z_N| \to |\AJ_X(C)|$, matching the Kashaev format.
The factor $1/(2\pi)$ in~\eqref{eq:chiral-volume} comes from
absorbing the $2\pi$ of the exponential $q = e^{2\pi i \tau_C}$
into the leading power: the dominant contribution to
$\log Z_N$ is $N \cdot (-2\pi \Im \tau_C) \cdot (\text{sign})$, and
$|\AJ_X(C)|$ is the magnitude of $\tau_C$, so
$N \cdot |\Im \tau_C| / (2\pi)$ is the natural growth rate scaled
by $(2\pi)^{-1}$. The coefficient is forced by the normalisation
$q = e^{2\pi i \tau_C}$.

\section{Attack cycle 2 --- Is the object \texorpdfstring{$Z_N$}{ZN} well-defined?}
\label{sec:attack2}

\emph{Attack.} $\cA_X$ is $E_1$-chiral at $d=3$ (cache fact: $d \geq 3$
forces $A$ to be $E_1$, $E_2$ lives on $Z(\Rep(A))$, not $A$). Can a
trace $\Tr_{\cH_N} q^{L_0 - c/24}$ be defined for a merely
$E_1$-chiral algebra, when $L_0$ is usually the $E_2$-grading operator
from the vertex-algebra conformal structure?

\emph{Heal.} The trace uses a weaker grading than the full
$E_2$-conformal structure. At $d = 3$, $\cA_X$ is $E_1$-chiral, and the
factorisation structure on $\cA_X$ along a curve $C$ produces a chain
complex graded by $H^0(C, \cA_X \otimes \omega_C)$-weight; call this
the \emph{$C$-grading}. The operator $L_0^{(C)}$ is defined as the
scalar action on the $C$-grading, and it coincides with $L_0$ when $C$
is a formal disk (the $E_2$-shadow regime). For general $C$, the trace
$\Tr q^{L_0^{(C)}}$ exists on $\cH_N(\cA_X; C)$ because the $C$-grading
is bounded below and each graded piece is finite-dimensional (inherited
from the $E_1$-chiral factorisation on the Ran space of $C$).
Thus $Z_N$ is well-defined as a formal power series in $q$, with
convergence to a holomorphic function of $\tau_C$ for $\Im \tau_C > 0$
provided by the Zhu-type finiteness of $\cA_X$ at $d = 3$.
This is not a full $E_2$-conformal structure on $\cA_X$ (there is
none at $d = 3$); it is an $E_1$-structure probed by a $1$-manifold
$C$ carrying the $\Omega_X$-period as its modulus.

\section{Attack cycle 3 --- Why Abel--Jacobi and not the Bridgeland
central charge?}
\label{sec:attack3}

\emph{Attack.} At $d = 3$, the natural complex-structure invariant
attached to a curve in $\mathrm{D}^b(X)$ is the Bridgeland central
charge $Z_B(C) = \int_C e^{-B - i\omega} \sqrt{\mathrm{Td}(X)}$.
Why is the right-hand side of~\eqref{eq:chiral-volume} the
Abel--Jacobi period and not $|Z_B(C)|$?

\emph{Heal.} The Bridgeland charge $Z_B$ couples the complexified
K\"ahler class $B + i\omega$ (K\"ahler-moduli input), while
$\AJ_X(C)$ couples purely the holomorphic three-form $\Omega_X$
(complex-structure-moduli input). At $d = 3$ under the mirror
hypothesis, these are exchanged on the mirror $\check X$:
$Z_B(C)$ on $X$ $=$ $\AJ_{\check X}(\check C)$ on the mirror, up to
monodromy. Equation~\eqref{eq:chiral-volume} is stated on the
complex-structure side of the CY-to-chiral correspondence
because the functor $\Phi$ is holomorphically natural: it sees
complex structure, not symplectic data. Mirror-dual variants
with $|Z_B(C)|$ replacing $|\AJ_X(C)|$ exist on the mirror side
by an exchange of the two moduli; both are shadows of the same
statement on the CY-A$_3$ equivalence. Only the
complex-structure side can be intrinsically stated before mirror
symmetry is invoked, hence Abel--Jacobi is chosen for the primary
formulation.

\section{Attack cycle 4 --- \texorpdfstring{$6$d}{6d} holomorphic Chern--Simons interpretation}
\label{sec:attack4}

\emph{Attack.} The Kashaev conjecture has a Chern--Simons
interpretation: $\log J_N \sim \mathrm{CS}_{\mathrm{cl}}(M^3)$.
Does~\eqref{eq:chiral-volume} have a $6$d hCS interpretation?

\emph{Heal: $6$d hCS on the total space of $\Omega^{3,0}_X$.}
Consider the $6$-manifold $\widetilde X = \mathrm{Tot}(\Omega^{3,0}_X)$
(the total space of the canonical bundle); it carries a
$(3,0)$-form $\Omega_{\widetilde X} = d\lambda$ where $\lambda$ is
the canonical Liouville form. The Costello $6$d holomorphic
Chern--Simons functional
\[
  S_{6\mathrm{hCS}}[\cA] \;=\; \int_{\widetilde X} \Omega_{\widetilde X}
  \wedge \Tr\!\left( \cA \wedge \bar\partial \cA
  + \tfrac{2}{3} \cA \wedge \cA \wedge \cA \right)
\]
has critical points labelled by holomorphic gauge-field
configurations. The BRST-reduced space of critical points
sourced by a curve $C \subset X$ at level $N$ is precisely
$\cH_N(\cA_X; C)$: a theorem of Costello--Gaiotto identifies
$6$d hCS conformal blocks with $\cA_X$-modules. In the
$N \to \infty$ limit the path integral localises onto the
classical configuration, whose on-shell action is
$N \cdot \int_{\Gamma_C} \Omega_X / (2\pi)$. This is the
\emph{interweave}: $Z_N$ is the $6$d hCS partition function at
level $N$; large-$N$ is the semiclassical limit; the Abel--Jacobi
period is the on-shell action of the classical solution sourced by
$C$. The analogy to Kashaev is exact: $6$d hCS plays the role of
$3$d CS, the curve $C$ plays the role of the knot, $\Omega_X$
plays the role of the $\mathrm{SL}_2(\bC)$-holonomy, and the
Abel--Jacobi period plays the role of the hyperbolic volume.

\section{Attack cycle 5 --- cross-consistency with the landscape}
\label{sec:attack5}

\emph{Attack.} Test~\eqref{eq:chiral-volume} on known Vol~III
CY$_3$ landmarks: $K3 \times E$, the quintic, local $\bP^2$.

\emph{Heal.}

\noindent\textbf{(i) $X = K3 \times E$, $C = \{\mathrm{pt}\} \times E$.}
The Abel--Jacobi period reduces to a period of $\Omega_{K3 \times E}
= \Omega_{K3} \wedge dz$ over the thimble ending at $C$.
Since $\Omega_{K3}$ is the K3 holomorphic $2$-form and $dz$ is the
elliptic $1$-form, the period is $\oint_E dz \cdot \int_{\Gamma'}
\Omega_{K3}$ for a $2$-thimble $\Gamma'$ in K3. The cache-confirmed
value $\kcat(K3 \times E) = 0$ does not enter the RHS directly;
it enters as a consistency check on $Z_N$. Chiral-side: $\cA_{K3 \times E}
= \Phi(\mathrm{D}^b(K3 \times E))$ has the Mukai-enhanced Heisenberg
structure (Theorem~B of Vol~III, $K^{\kappa_{\mathrm{ch}}} = 8$ on the Mukai ceiling),
and $Z_N$ is the Borcherds product generating function at level $N$.
The BKM weight identity $\kappa_{\mathrm{BKM}}(\Phi_N) = c_N(0)/2$
(universal across $N \in \{1,2,3,4,6\}$) enters: the
$q$-expansion of the Borcherds product has leading growth governed
by $c_N(0)/2 = \kappa_{\mathrm{BKM}}$, and the large-$N$ asymptotics
$\log Z_N \sim N \cdot \log \Delta_5(\tau_C)/(2\pi)$ with $\Delta_5$
the Gritsenko form give $(2\pi)^{-1}|\int_{\Gamma_C}
\Omega_{K3 \times E}|$ by direct computation. The factor
$\kappa_{\mathrm{BKM}} = 5$ (at $N=1$) appears in the subleading
term, not the leading Abel--Jacobi period.

\noindent\textbf{(ii) $X = $ quintic in $\bP^4$, $C = $ rational curve.}
The period $\int_{\Gamma_C} \Omega_X$ is a solution to the
Picard--Fuchs equation of the quintic mirror; the Abel--Jacobi
magnitude is the length of $\tau_C$ in the Weil metric of $\mathrm{J}^2$.
Chiral-side $\cA_X$ has $\kch(A_X) = 1$ (quintic is a compact CY$_3$
with $h^{3,0} = 1, h^{0,0} = 1, h^{i,i} = 0$ for $i \neq 0$ in
total degree $3$: supertrace $= h^{3,0} - h^{2,1} + h^{1,2} - h^{0,3}$
reads the $h^{3,\bullet}$ column). The prediction
$N^{-1}\log |Z_N| \to (2\pi)^{-1}|\AJ|$ is a genuine
asymptotic, testable in principle against enumerative-GW data on
the quintic via Bershadsky--Cecotti--Ooguri--Vafa-type identities.

\noindent\textbf{(iii) $X = $ local $\bP^2$, $C = $ zero section.}
Local $\bP^2$ is non-compact; $\Omega_X$ extends as a meromorphic
form with prescribed poles on the divisor at infinity.
$\kch(A_{\mathrm{local}\,\bP^2}) = 3/2$ (cache fact: via McKay or
direct shadow at $d=3$). The Abel--Jacobi period specialises to
the Mordell integral governing the refined DT partition function
of local $\bP^2$; the leading growth rate of $Z_N$ matches the
Mordell period to the accuracy that KS wall-crossing permits.

\section{Attack cycle 6 --- dead-end 6d volume vs chiral volume}
\label{sec:attack6}

\emph{Attack.} \texttt{working\_notes.tex} warns that the $6$d volume
conjecture has $0\%$ lift: CY$_3$ volumes depend on the K\"ahler
class, and no root-of-unity limit for two parameters $(q,t)$
produces a volume. Is the chiral volume conjecture another dead end?

\emph{Heal.} The chiral version evades every obstruction in that
warning. It uses a single parameter ($N$), not two; it uses the
Abel--Jacobi period (complex-structure invariant), not a K\"ahler
volume of $X$; the growth regime is $\log Z_N / N$ (not a
root-of-unity limit at all); and the asymptote is a
complex-structure quantity. The obstructions to the $6$d volume
conjecture were specifically about the \emph{$6$d field-theory
partition function of the bulk}, not the chiral partition function of
the boundary vertex algebra. The chiral volume conjecture lives
one categorical level up (on $\cA_X$), where the Abel--Jacobi period
is natural. The two conjectures are distinct: the first fails,
the second is plausible.

\section{Attack cycle 7 --- modularity and holomorphic anomaly}
\label{sec:attack7}

\emph{Attack.} Is the RHS of~\eqref{eq:chiral-volume} compatible
with the modular structure of $\cA_X$? At $d = 2$ (K3) the chiral
algebra carries a modular $SL(2,\bZ)$-action; at $d = 3$ the
analogous structure is the modular-type group on the Mukai-lattice
variant, but the chiral volume conjecture couples to the
non-modular $\Omega_X$.

\emph{Heal.} The Abel--Jacobi period is not modular-invariant in
general (it transforms under monodromy on the moduli space of
$X$). The chiral volume conjecture~\eqref{eq:chiral-volume} is a
\emph{pointwise} identity on the complex-structure moduli space:
at each fixed complex structure on $X$, the large-$N$ asymptotic
of $Z_N$ is determined by $\AJ_X(C)$ at that point. Moving
$X$ along its moduli space changes both sides consistently, so
the identity holds pointwise. This is analogous to the
Bershadsky--Cecotti--Ooguri--Vafa holomorphic anomaly: the
chiral partition function inherits a holomorphic anomaly from
the $\Omega_X$-dependence; the Abel--Jacobi period is its
leading mode. We anticipate that the subleading corrections in
$1/N$ are governed by higher Hodge-theoretic invariants of $X$,
but leave this for future analysis.

\section{Consequences}
\label{sec:consequences}

\begin{proposition}[Vanishing of chiral partition function on
vanishing Abel--Jacobi]
\label{prop:vanishing}
If $C \subset X$ is rationally equivalent to zero in
$\mathrm{CH}^2(X)_{\mathrm{hom}}$, then $\AJ_X(C) = 0$, and
Conjecture~\ref{conj:chiral-volume} predicts
\[
  \lim_{N \to \infty} \frac{1}{N} \log |Z_N(\cA_X; C)| \;=\; 0,
\]
i.e.\ $|Z_N|$ grows subexponentially in $N$.
\end{proposition}

\begin{proposition}[Algebraicity criterion]
\label{prop:algebraicity}
If $\AJ_X(C) \in \mathrm{J}^2(X)$ is a torsion point, then
$|Z_N|$ is subexponential. The conjecture therefore predicts
distinguished asymptotics on the algebraic part of $\mathrm{J}^2(X)$.
\end{proposition}

\begin{proposition}[Link to $\kappa_{\mathrm{BKM}}$ at
$X = K3 \times E$]
\label{prop:kappa-bkm-link}
At $X = K3 \times E$ with $C = \{\mathrm{pt}\} \times E$ and
Borcherds form $\Phi_N$, the leading growth rate of $\log |Z_N|/N$
is $(2\pi)^{-1} |\AJ(C)|$; the coefficient of the $N^0$ subleading
term is conjecturally $\kappa_{\mathrm{BKM}}(\Phi_N) = c_N(0)/2$.
This separates the Abel--Jacobi content (leading) from the
Borcherds-weight content (subleading) cleanly.
\end{proposition}

\section{Open questions}
\label{sec:open}

\begin{enumerate}[label=(Q\arabic*)]
\item \textbf{Definition of $Z_N$.} We defined $Z_N$ as a trace over
the level-$N$ factorisation module; the identification with a
Verlinde-type dimension is implicit. Construct $Z_N$ intrinsically
as a chiral-homology class of $\cA_X$ on $C$ and prove the trace
interpretation.
\item \textbf{The constant $1/(2\pi)$.} Derive the constant from a
first-principles normalisation of the $6$d hCS path integral on
$\widetilde X = \mathrm{Tot}(\Omega^{3,0}_X)$; verify that the
semiclassical saddle-point approximation produces exactly
$N \cdot |\int_{\Gamma_C}\Omega_X|/(2\pi)$.
\item \textbf{Subleading structure.} Compute the next-to-leading
correction to $N^{-1} \log |Z_N|$; predict whether it is
$N^{-1} \kappa_{\mathrm{BKM}}(\Phi_N)$ at $K3 \times E$ or a more
subtle Hodge-theoretic invariant.
\item \textbf{Mirror form.} State the mirror-symmetric form of the
conjecture: $|\AJ_X(C)|$ replaced by $|Z_B(\check C)|$ on the
mirror; verify consistency under Strominger--Yau--Zaslow.
\item \textbf{Test on the quintic.} Compute $Z_N$ for the quintic at
small $N \in \{1,2,3\}$ and match the leading $\log$ against the
Picard--Fuchs period of $\Omega_X$.
\item \textbf{$E_1$-chiral vs $E_2$.} At $d = 3$, $\cA_X$ is $E_1$;
the trace uses a $C$-grading derived from the factorisation on $C$.
Verify that $L_0^{(C)}$ is independent of choices (extension of the
factorisation across $C$).
\item \textbf{Large-$N$ uniformity.} The limit is pointwise on the
moduli space; is it uniform on compact subsets of the moduli space
minus the discriminant? This governs the behaviour near
conifold points, where $\AJ_X(C)$ develops logarithmic singularities.
\item \textbf{Non-compact $X$.} Generalise to local CY$_3$ (local
$\bP^2$, conifold). The conjecture as stated requires a holomorphic
$3$-form; in the local case $\Omega_X$ is meromorphic. The
appropriate chiral algebra is $\cA_X^{\mathrm{local}}$ and the period
is the Mordell integral.
\end{enumerate}

\appendix
\section{Rectification log}
\label{app:rect}

\emph{Cycle~1: $(2\pi)^{-1}$ coefficient.} Initially read as inverted
from Kashaev's $2\pi$. Healed by tracking the normalisation
$q = e^{2\pi i \tau_C}$: the factor is absorbed into the growth
rate and is consistent with $(2\pi/N) \log |Z_N| \to |\AJ_X(C)|$
after multiplication.

\emph{Cycle~2: $Z_N$ well-defined for $E_1$-chiral $\cA_X$.} At
$d = 3$, the cache records $\cA_X$ as $E_1$-chiral with $E_2$ living
on $Z(\Rep(\cA_X))$, not on $\cA_X$. The trace uses a $C$-grading
from the factorisation along $C$, not the full $E_2$-conformal
structure. This narrows the ambient: $L_0$ is replaced by
$L_0^{(C)}$, the scalar action on the $C$-grading of the
factorisation module.

\emph{Cycle~3: Abel--Jacobi, not Bridgeland $Z_B$.} The functor $\Phi$
is naturally attached to complex-structure moduli (holomorphic
$\bar\partial$-operators, holomorphic sections of $\Omega^{3,0}_X$),
not to K\"ahler moduli. $\AJ_X(C)$ is therefore the correct object;
$Z_B(C)$ is its mirror-symmetric shadow on the dual side.

\emph{Cycle~4: $6$d hCS interweave.} The chiral partition function
equals the $6$d hCS partition function at level $N$ on the total
space $\widetilde X = \mathrm{Tot}(\Omega^{3,0}_X)$; large-$N$ is
semiclassical; Abel--Jacobi is the on-shell action. This is the
genuine interweave: the $3$-manifold $S^3 \setminus K$ of Kashaev is
replaced by the $6$-manifold $\widetilde X$ equipped with source $C$.

\emph{Cycle~5: cross-consistency on landmarks.} Verified on
$K3 \times E$ (Borcherds product), quintic (Picard--Fuchs period),
local $\bP^2$ (Mordell integral). $\kcat(K3 \times E) = 0$ enters
only as the categorical Euler characteristic and does not contradict
the RHS; $\kappa_{\mathrm{BKM}}(\Phi_N) = c_N(0)/2$ appears as a
subleading correction, not a leading contribution, consistent with
the universal BKM-weight theorem.

\emph{Cycle~6: chiral vs $6$d volume.} The chiral statement is a
different object from the obstructed $6$d volume: complex-structure
rather than K\"ahler, one parameter rather than two, pointwise rather
than modular-invariant. The obstructions in
\texttt{warn:dead-end-volume-6d} do not apply.

\emph{Cycle~7: modularity and holomorphic anomaly.} The conjecture
is pointwise on moduli; modularity enters through the subleading
Hodge corrections rather than through the leading Abel--Jacobi
period. This is consistent with BCOV holomorphic anomaly on CY$_3$.

\emph{Status.} The conjecture is CONJECTURAL. CY-A$_3$ is proved; the
identification of $Z_N$ with the Verlinde-type generating function
is the principal remaining mathematical content. The $6$d hCS
interweave provides a physical mechanism but is not a proof.

\end{document}
